\section{SmallWood Proof-System Details}
\label{app:smallwood}

This appendix recalls the SmallWood proof stack and provides algorithmic details used by our instantiations. The main
text treats SmallWood as a black box; here we summarize the DECS/LVCS/PCS/PIOP layers and the Fiat--Shamir compilation,
following \cite{EPRINT:FenRiv25b}.

\newcommand{\NTT}{\mathrm{NTT}}

\subsection{Domains and packing}
Fix a base field \(\Fq\), an evaluation set \(\Omega\subset\Fq\) of size \(s\), and a disjoint tail set
\(\Omega'\subset\Fq\setminus\Omega\) of size \(\ell\). A witness row is interpolated into a polynomial of degree
\(\le s+\ell-1\) whose values on \(\Omega\) are the packed coordinates and whose values on \(\Omega'\) are random masks
(HVZK).

\subsection{DECS (Degree-Enforcing Commitment Scheme)}
DECS commits to evaluations of witness polynomials and enforces that they come from low-degree polynomials rather than
arbitrary vectors. Commitments are implemented via a Merkle tree over an NTT-sized domain, and degree soundness is
obtained by checking masked low-degree relations at verifier-chosen indices.

\begin{algorithm}[H]
\caption{\textsf{DECS.DeriveGamma}(\(\mathsf{root},\eta,r,q\)) \(\to \Gamma\) (verifier and prover)}
\label{alg:decs-derivegamma}
\begin{algorithmic}[1]
\State Initialize counter \(ctr\gets 0\); let \(L\gets \big\lfloor \frac{2^{64}}{q}\big\rfloor\cdot q\).
\For{\(k=0\) to \(\eta-1\)}
  \For{\(j=0\) to \(r-1\)}
    \Repeat
      \State \(x \gets \textsf{SHA256}(\mathsf{root}\ \|\ \textsf{LE64}(ctr))\) interpreted as a 64-bit integer
      \State \(ctr\gets ctr+1\)
    \Until{\(x < L\)}
    \State \(\Gamma_{k,j}\gets x \bmod q\)
  \EndFor
\EndFor
\State \Return \(\Gamma\)
\end{algorithmic}
\end{algorithm}

\begin{algorithm}[H]
\caption{\textsf{DECS.VerifyEval}(\(\mathsf{root},\Gamma,R,\mathsf{open}\)) (verifier)}
\label{alg:decs-verifyeval}
\begin{algorithmic}[1]
\Require Domain size \(N\), modulus \(q\), mask count \(\eta\), committed row count \(r\).
\For{each opened index \(e\)}
  \State Verify Merkle authentication for \(e\) against \(\mathsf{root}\); reject on failure.
  \For{each mask row \(k<\eta\)}
    \State Check \(R_k^\star[e] \stackrel{?}{=} M_k(e) + \sum_{j=0}^{r-1}\Gamma_{k,j}\,P_j(e)\pmod q\).
  \EndFor
\EndFor
\State \Return accept
\end{algorithmic}
\end{algorithm}

\subsection{LVCS (Linear Vector Commitments)}
LVCS commits to packed witness rows and supports succinct openings of \(\Fq\)-linear forms of those rows. SmallWood uses
LVCS both as a routing/batching layer and as the core of its hash-based PCS: the verifier requests openings of linear
forms that correspond to evaluating constraint residuals at verifier-chosen points, and DECS enforces that opened values
come from low-degree polynomials.

\begin{algorithm}[H]
\caption{\textsf{LVCS.CommitInitWithParams}(\(\mathrm{ringQ},\{r_j\}_{j<r},\ell,\mathrm{params}\)) \(\to (\mathsf{root},\mathsf{ProverKey})\) (prover)}
\label{alg:lvcs-commitinit}
\begin{algorithmic}[1]
\Require Ring \(R_q\) of size \(N\), packed rows \(r_j\in\Fq^{n_{\mathrm{cols}}}\), tail length \(\ell>0\).
\State Sample \(\mathrm{mask}_j\getsunif \Fq^{\ell}\) for each row \(j\) (HVZK tails).
\State Interpolate each row to \(P_j(X)\gets \textsf{interpolateRow}(\mathrm{ringQ},r_j,\mathrm{mask}_j,n_{\mathrm{cols}},\ell)\).
\State Commit to \(\{P_j\}\) via DECS and obtain Merkle root \(\mathsf{root}\).
\State Derive \(\Gamma\gets \textsf{DECS.DeriveGamma}(\mathsf{root},\eta,r,q)\) and cache values needed for openings.
\State \Return \((\mathsf{root},\mathsf{ProverKey})\).
\end{algorithmic}
\end{algorithm}

\begin{algorithm}[H]
\caption{\textsf{LVCS.CommitStep1}(\(\mathsf{root}\)) \(\to \Gamma\) (verifier)}
\label{alg:lvcs-commitstep1}
\begin{algorithmic}[1]
\State Store \(\mathsf{root}\) and derive \(\Gamma\gets \textsf{DECS.DeriveGamma}(\mathsf{root},\eta,r,q)\).
\State \Return \(\Gamma\).
\end{algorithmic}
\end{algorithm}

\begin{algorithm}[H]
\caption{\textsf{LVCS.CommitStep2}(\(\{R_k\}_{k<\eta}\)) \(\to\) accept/reject (verifier)}
\label{alg:lvcs-commitstep2}
\begin{algorithmic}[1]
\For{each mask row \(R_k\)}
  \State Reject if \(\deg(R_k) > d\) (degree guard before openings).
\EndFor
\State \Return accept.
\end{algorithmic}
\end{algorithm}

\begin{algorithm}[H]
\caption{\textsf{LVCS.EvalInitMany}(\(\mathrm{ringQ},\mathsf{ProverKey},\{\mathrm{Coeffs}_k\}_{k<m}\)) \(\to \{\bar v_k\}_{k<m}\) (prover)}
\label{alg:lvcs-evalinit}
\begin{algorithmic}[1]
\Require Linear forms \(\mathrm{Coeffs}_k\in\Fq^{r}\) for \(k=0,\dots,m-1\).
\State Let \(\ell\gets |\mathrm{mask}_j|\).
\For{\(k=0\) to \(m-1\)}
  \State Initialize \(\bar v_k\gets 0\in\Fq^{\ell}\).
  \State Set \(\bar v_k \gets \sum_{j=0}^{r-1} \mathrm{Coeffs}_k[j]\cdot \mathrm{mask}_j \in \Fq^{\ell}\).
\EndFor
\State \Return \(\{\bar v_k\}_{k<m}\).
\end{algorithmic}
\end{algorithm}

\begin{algorithm}[H]
\caption{\textsf{LVCS.ChooseE}(\(\ell,n_{\mathrm{cols}}\)) \(\to E\) (verifier)}
\label{alg:lvcs-choosee}
\begin{algorithmic}[1]
\Require Ring size \(N\); packed domain size \(n_{\mathrm{cols}}=|\Omega|\); masked slab indices \([n_{\mathrm{cols}},n_{\mathrm{cols}}+\ell)\).
\State Sample \(E\subset \{n_{\mathrm{cols}}+\ell,\dots,N-1\}\) uniformly without replacement with \(|E|=\ell\).
\State \Return \(E\).
\end{algorithmic}
\end{algorithm}

\begin{algorithm}[H]
\caption{\textsf{LVCS.EvalFinish}(\(\mathsf{ProverKey},E\)) \(\to \mathsf{Opening}\) (prover)}
\label{alg:lvcs-evalfinish}
\begin{algorithmic}[1]
\State \(\mathsf{open}\gets \mathsf{ProverKey}.\mathrm{DecsProver}.\textsf{EvalOpen}(E)\).
\State \Return \(\mathsf{Opening}\{\mathrm{DECSOpen}=\mathsf{open}\}\).
\end{algorithmic}
\end{algorithm}

\begin{algorithm}[H]
\caption{\textsf{LVCS.EvalStep2}(\(\{\bar v_k\},E,\mathsf{open},C,v_{\mathrm{targets}}\)) (verifier)}
\label{alg:lvcs-evalstep2}
\begin{algorithmic}[1]
\Require Coeffs \(C\in\Fq^{m\times r}\), public prefixes \(v_{\mathrm{targets}}\in(\Fq^{n_{\mathrm{cols}}})^{m}\).
\State Split \(\mathsf{open}\) into mask openings for \([n_{\mathrm{cols}},n_{\mathrm{cols}}+\ell)\) and tail openings at \(E\).
\State Verify the DECS openings (degree soundness and Merkle authentication); reject on failure.
\For{each \(k<m\) and each masked index \(i<\ell\)}
  \State Check \(\sum_{j=0}^{r-1} C[k,j]\cdot \mathsf{maskOpen.Pvals}[i][j] \stackrel{?}{=} \bar v_k[i]\pmod q\).
\EndFor
\State Reconstruct each batched polynomial \(Q_k(X)\) from its \(\Omega\)-prefix \(v_{\mathrm{targets}}[k]\) and mask tail \(\bar v_k\).
\State Check tail consistency of \(Q_k\) against the opened linear combinations at points in \(E\).
\State \Return accept.
\end{algorithmic}
\end{algorithm}

\subsection{LVCS and PCS/PIOP}
LVCS provides linear opening queries to committed rows and is used to implement the polynomial oracle interface needed
by PIOPs. SmallWood combines LVCS with DECS (for degree enforcement) to obtain a hash-based PCS.

\subsection{What DECS/LVCS/PCS bind in ARC-SPRUCE}\label{app:smallwood-binds}
At a high level, the SmallWood stack provides a single commitment that simultaneously binds \emph{all} witness rows used
in a statement, and it enforces that every opened value is consistent with a low-degree polynomial view of those rows.
In ARC-SPRUCE this binding is used in three complementary ways:
\begin{itemize}[leftmargin=1.25em]
  \item \textbf{Row binding and low degree (DECS).} DECS ensures that the prover cannot answer openings using
  high-degree ``fake'' vectors: each committed row must correspond to a degree-\(\le s+\ell-1\) polynomial that matches
  the packed values on \(\Omega\) and the masked tails on \(\Omega'\).
  \item \textbf{Linear linking across subsystems (LVCS).} LVCS is used to request openings of verifier-chosen linear
  combinations of rows. This is what couples different parts of the statement: the same opened evaluations are reused
  to recompute commitment residuals, centering residuals, the rational-hash residual, the signature residual, and the
  PRF checkpoint residuals.
  \item \textbf{No witness splitting (single PCS commitment).} Because all rows live under one commitment, the prover
  cannot use inconsistent witnesses for different checks. For instance, in showings the same signed packed-message rows
  that appear in the hash/signature relation are also the rows from which the committed PRF key lanes are extracted;
  this prevents ``mix-and-match'' of \((k,u)\) with an unrelated PRF witness.
\end{itemize}
This is exactly the binding invariant emphasized in the previous paper’s SmallWood appendix: the verifier’s checks are
expressed as residual functions of opened row evaluations, so all constraints are evaluated against a single coherent
row assignment.

\subsection{Openings we request in issuance vs showing}\label{app:smallwood-openings}
Both proofs in ARC-SPRUCE follow the same pattern: the prover commits to all witness rows as polynomials, the verifier
derives Fiat--Shamir challenges, and the prover opens the committed rows at verifier-chosen points. The verifier then
recomputes residuals for each constraint family from these openings and checks that all residuals vanish.

\paragraph{Issuance (pre-sign).}
The witness consists of the base rows
\[
  M_1,M_2,R_0^U,R_1^U,\bar R,\quad R_0,R_1,K_0,K_1.
\]
while the target \(T\) and issuer challenges \((r_0^I,r_1^I)\) are public. The verifier requests openings sufficient to
replay: (i) the linear commitment constraints linking \(\mathsf{com}\) to \((M_1,M_2,R_0^U,R_1^U,\bar R)\), (ii) the
centering constraints linking \((R_0,R_1)\) to the holder/issuer randomness and carries \((K_0,K_1)\), and (iii) the
cleared-denominator hash constraint linking \(T\) to \((M_1,M_2,R_0,R_1)\), together with packing and membership checks.
Because all of these residuals are evaluated from the same opened rows, the proof binds the target \(T\) to the
commitment opening and the centered randomness.

\paragraph{Showing (post-sign).}
The witness extends issuance with the hash target \(T\), the signature preimage rows \(U\), and the PRF witness block
(key lanes and checkpointed S-box outputs). The verifier requests openings sufficient to replay: (i) the signature
residual \(A\cdot U-T\), (ii) the bilinear cleared-denominator hash residual linking \(T\) to the base rows, and (iii)
the PRF key-binding, checkpoint, and tag-binding residuals (Section~\ref{subsec:prf-checkpointed}). The key point is
that the PRF key lanes are tied back to the signed packed message via selector constraints, so the opened values jointly
bind \(\mathsf{tag}=\PRF(k,\nonce)\) to the same secret \(k\) that appears inside the signed message
\(m=\mu\Vert k\).

\paragraph{Internal \(\Omega\)-sum check.}
Even when an application’s aggregated family is empty (\(F_{\mathrm{agg}}=\emptyset\)), PACS/PIOP still enforces the
masked \(\Omega\)-sum test on the batch polynomials \(Q\), which is the mechanism that ties the verifier’s spot-check
openings to correctness on the full packing domain \(\Omega\).

\subsection{PACS/PIOP interface (used by our statements)}
SmallWood’s PACS/PIOP layer checks that constraint polynomials vanish on \(\Omega\) (and at random points outside
\(\Omega\)), by committing to witness rows and mask polynomials and then opening them at verifier-chosen coordinates.
The following algorithms summarize the high-level structure used in our proofs.

\begin{algorithm}[H]
\caption{\textsf{PACS.PIOPCommitOracle}\(_{\mathrm{Prover}}(\textsf{params})\to \textsf{pcsCom}\)}
\label{alg:pacs-commitoracle}
\begin{algorithmic}[1]
\State Interpolate witness rows into \(P_1,\ldots,P_n\in\Fq[X]_{\le s+\ell'-1}\) with \(\ell'\) random blind evaluations outside \(\Omega\).
\State Sample mask polynomials \(M_1,\ldots,M_\rho\in\Fq[X]_{\le d_Q}\) such that \(\sum_{\omega\in\Omega} M_i(\omega)=0\) for all \(i\).
\State Commit to \(P\Vert M\) via PCS (implemented through LVCS/DECS); obtain \(\textsf{pcsCom}\).
\State \Return \(\textsf{pcsCom}\).
\end{algorithmic}
\end{algorithm}

\begin{algorithm}[H]
\caption{\textsf{PACS.PIOPBatchCombine}\(_{\mathrm{Prover}}(\textsf{pcsCom},\Gamma')\to Q\)}
\label{alg:pacs-batchcombine}
\begin{algorithmic}[1]
\State Using \(P,M\), form batched polynomials
\[
Q_i(X) \;=\; M_i(X) + \sum_j \Gamma'_{i,j}(X)\,F_j(X) + \sum_u \gamma'_{i,u}\,F'_u(X),
\]
where \(F_j\) and \(F'_u\) are the parallel and aggregated constraint polynomials evaluated on \(P\).
\State \Return \(Q=(Q_1,\ldots,Q_\rho)\).
\end{algorithmic}
\end{algorithm}

\begin{algorithm}[H]
\caption{\textsf{PACS.PIOPOpenAndCheck}(\(\textsf{pcsCom},Q\)) (verifier)}
\label{alg:pacs-openandcheck}
\begin{algorithmic}[1]
\State Sample a query set \(E'\subset S\) with \(|E'|=\ell'\) (Fiat--Shamir in NIZK).
\State Verify PCS openings of \(P\) and \(M\) at \(E'\) (via LVCS/DECS); reject on failure.
\State Check that each \(Q_i(e)\) matches its definition from \(M_i(e)\) and the constraint evaluations at \(e\in E'\).
\State Check the masked \(\Omega\)-sum test \(\sum_{\omega\in\Omega} Q_i(\omega)=0\) for all \(i\).
\State \Return accept/reject.
\end{algorithmic}
\end{algorithm}

\subsection{Fiat--Shamir compilation}
SmallWood compiles the interactive protocol to a NIZK using Fiat--Shamir with grinding to obtain challenges with
all-zero prefixes, as in \cite{EPRINT:FenRiv25b}. Our ARC instantiation is dominated by parallel constraints, but
the PACS/PIOP reminder here is stated for general \((F_{\mathrm{par}},F_{\mathrm{agg}})\) statements. In particular, if the
application introduces global coupling across the packing domain (as in the proof-friendly encodings inherited from the
previous paper), those relations appear as \(F_{\mathrm{agg}}\) and are enforced through the \(\Omega\)-sum checks of the PIOP.

\begin{algorithm}[H]
\caption{\textsf{PACS.ToNIZK} (Fiat--Shamir wrapper with grinding)}
\label{alg:pacs-tonizk}
\begin{algorithmic}[1]
\State Sample \(\textsf{salt}\getsunif\{0,1\}^{2\lambda}\).
\State Derive challenges by hashing the transcript with grinding counters to obtain required prefix properties.
\State Output PCS/LVCS/DECS openings and any required batched values; the verifier recomputes challenges and checks.
\end{algorithmic}
\end{algorithm}
